% !TeX root = ../../Book2.tex
% 纲要标题：### 14.1 Jacobson 根
% 纲要位置：计划纲要.md，第 1492 行。

\section{Jacobson 根}
\label{b2:sec:1401}

简单模给出环最小的非零作用对象。某个环元素可能在每个简单模上都作用为零；这些元素组成 Jacobson 根。本节把这个描述与极大左理想、可逆性及取商联系起来。环 \(R\) 均含单位元，模均幺性，不假设 \(R\) 交换或满足链条件。

% 纲要要点（供目录核对）：
%
% * 极大左理想的交及简单模作用的共同核；极大左理想商首先为模
% * 根中的可逆性判据及 Jacobson 逆公式；区分单个幂零元与理想中元素各自幂零
% * 商环、矩阵环与有限长度判据；矩阵环简单模对应、单位提升及上三角根计算
%

\subsection{极大左理想的交及等价刻画}
\label{b2:subsec:1401:01}

每个简单左模都同构于某个 \(R/L\)，其中 \(L\) 为极大左理想。\cref{b2:prop:simple-annihilator-intersection}已经证明：一个元素零化全部简单模，当且仅当它属于全部极大左理想。因而这些左理想的交正好把简单作用共同检测不到的部分收集起来。这里 \(R/L\) 首先是简单左模；\(L\) 未必为双边理想，不能一般地把它当作商环。

\begin{definition}\label{b2:def:jacobson-radical}
设 \(R\) 为含幺环。将 \(R\) 的全部极大左理想取交，所得集合称为 \(R\) 的\term[Jacobson-gen]{Jacobson 根}[Jacobson radical]，记为
\[
J(R)=\bigcap_{L\text{ 极大左理想}}L.
\]
对零环约定 \(J(0)=0\)。非零环的每个真左理想都包含于某个极大左理想。
\end{definition}
\symidx{\(J(R)\)：环 \(R\) 的 Jacobson 根}

最后的存在性可由第一卷附录A定理A.1.4的 Zorn 引理得到：包含给定真左理想的真左理想族非空，链的并仍是左理想，而且不含一，因而仍真；空链用原理想作上界。由此取到极大元。这里“极大”在所有真左理想之间比较，不要求它本身是双边理想。

\begin{proposition}\label{b2:prop:radical-simple-annihilators}
设 \(R\) 为含幺环，\(J(R)\) 为其 Jacobson 根。\(J(R)\) 等于全部幺性简单左 \(R\) 模的零化子的交，因而是双边理想。对每个这样的简单模 \(S\)，令其作用同态为
\[
\rho_S:R\longrightarrow\operatorname{End}_{\mathbb Z}(S),
\qquad \rho_S(r)(s)=rs,
\]
其中自同态取自 \(S\) 的加法群，则 \(J(R)=\bigcap_S\ker\rho_S\)。两个交都可只取简单模的同构类代表；取全部循环商 \(R/L\)、\(L\) 为极大左理想，也足以计算它们。
\end{proposition}
\begin{proof}
交的等式已在\cref{b2:prop:simple-annihilator-intersection}证明：一个元素属于所有极大左理想，恰好使它在每个简单模的每个向量上作用为零。每个简单模本来就同构于某个循环商 \(R/L\)，同构模又有相同的零化子，所以只需所述代表族。模零化子是双边理想，因为 \(xS=0\) 蕴含 \((rx)s=r(xs)=0\)、\((xr)s=x(rs)=0\)。因此这些零化子的交 \(J(R)\) 也为双边理想。作用的结合律使 \(\rho_S\) 保持乘法，且 \(\ker\rho_S=\operatorname{Ann}_R(S)\)，便得到共同核的表述。零环没有非零幺性简单模，空族的交按全环取值，等式也成立。
\end{proof}

例如，对域 \(k\) 和 \(n\ge1\)，\(R=M_n(k)\) 的列向量模 \(k^n\) 简单且忠实，所以 \(J(R)=0\)。但是 \(x=E_{12}\) 在 \(n\ge2\) 时是非零幂零元。取 \(r=E_{21}\)，就有 \(rx=E_{22}\)，它是非零幂等元，已经不再幂零。单个元素幂零，并不保证任意左倍仍幂零；这与整个理想中的元素都幂零是不同的条件。

\subsection{根中的可逆性判据}
\label{b2:subsec:1401:02}

\begin{theorem}[可逆性判据]\label{b2:thm:jacobson-unit-criterion}
设 \(R\) 为含幺环，\(J(R)\) 为其 Jacobson 根。对 \(x\in R\)，以下条件等价：
\begin{equivalences}
\item\label{b2:item:radical-member} \(x\in J(R)\)。
\item\label{b2:item:radical-unit-left} 对每个 \(r\in R\)，\(1-rx\) 为单位。
\item\label{b2:item:radical-unit-right} 对每个 \(r\in R\)，\(1-xr\) 为单位。
\end{equivalences}
此外，对任意 \(a,b\in R\)，\(1-ab\) 可逆当且仅当 \(1-ba\) 可逆；若 \(u=(1-ab)^{-1}\)，则
\[
(1-ba)^{-1}=1+bua.
\]
\end{theorem}
\begin{proof}
先证明附带的恒等事实。若 \(u=(1-ab)^{-1}\)，则 \((1-ab)u=u(1-ab)=1\) 分别给出 \(u-1-abu=0\) 与 \(u-1-uab=0\)。于是
\[
(1-ba)(1+bua)=1+b(u-1-abu)a=1,
\qquad
(1+bua)(1-ba)=1+b(u-1-uab)a=1.
\]
所以 \(1+bua\) 是 \(1-ba\) 的双侧逆；交换 \(a,b\) 便得另一方向。

对根中的元素，要先由极大左理想的定义得到左逆，再利用该左逆仍具有 \(1+J(R)\) 的形式补出右逆。设 \(j\in J(R)\)。若左理想 \(R(1-j)\) 为真，它包含在某个极大左理想 \(L\) 中，而 \(j\in L\)，导致 \(1\in L\)，矛盾。故存在 \(b\) 满足 \(b(1-j)=1\)。此时 \(b=1+bj\in1+J(R)\)。对 \(-bj\in J(R)\) 重复同一论证，有 \(c\) 使 \(cb=1\)。于是
\[
c=cb(1-j)=1-j,
\]
所以 \((1-j)b=1\)，\(1-j\) 确实可逆。若 \(x\in J(R)\)，双边理想性给出 \(rx\in J(R)\)，因而得到第二项。

若第二项成立而 \(x\notin J(R)\)，取极大左理想 \(L\) 不含 \(x\)。此时 \(L+Rx=R\)，存在 \(\ell\in L,r\in R\) 使 \(1=\ell+rx\)。所以单位 \(1-rx=\ell\) 落在真左理想中，左乘其逆即得 \(1\in L\)，矛盾。

最后将已证恒等事实用于 \(a=r,b=x\)，便得第二、三项等价。
\end{proof}

其中 \(1-ab\) 与 \(1-ba\) 的可逆性等价及逆公式通常称为 Jacobson 引理。单位判据要求 \(1-rx\) 对任意 \(r\) 都可逆，描述的是所有左倍 \(rx\) 都不破坏 \(1\) 的可逆性；它为模掉根后的单位提升提供了依据。

\begin{corollary}[根的左右对称性]\label{b2:cor:jacobson-left-right}
对任意含幺环 \(R\)，有
\[
J(R)=\bigcap_{L\text{ 极大左理想}}L
=\bigcap_{K\text{ 极大右理想}}K.
\]
\end{corollary}
\begin{proof}
极大右理想正是反环 \(R^{\mathrm{op}}\) 的极大左理想。把\cref{b2:thm:jacobson-unit-criterion}用于 \(R^{\mathrm{op}}\)，可知 \(x\) 属于全部极大右理想，当且仅当对每个 \(r\in R\)，元素 \(1-xr\) 在 \(R\) 中可逆；反环与原环的单位相同。再把同一判据用于 \(R\)，这恰是 \(x\in J(R)\) 的条件。零环按约定也满足等式。
\end{proof}

单位判据中的“对每个 \(r\)”不能省去。前面的矩阵 \(x=E_{12}\) 满足 \((1-x)^{-1}=1+x\)，但取 \(r=E_{21}\) 后，\(1-rx=1-E_{22}\) 的第二列为零，不可逆。这正是单个幂零元不必属于根的障碍。

\begin{corollary}\label{b2:cor:nil-ideal-radical}
设 \(R\) 为含幺环、\(I\subseteq R\) 为双边理想。若 \(I\) 的每个元素都幂零，则 \(I\subseteq J(R)\)。特别地，若 \(I^n=0\) 对某个正整数 \(n\) 成立，结论也成立。
\end{corollary}
\begin{proof}
对 \(x\in I,r\in R\)，理想的左乘封闭性保证 \(rx\in I\)，所以 \(rx\) 也幂零；这是单个幂零元的例子中缺少的保证。取正整数 \(m\) 使 \((rx)^m=0\)，有限几何级数给出
\[
(1-rx)^{-1}=1+rx+\cdots+(rx)^{m-1}.
\]
由可逆性判据，\(x\in J(R)\)。理想幂零时，每个元素的同一次幂都为零，故为特例。
\end{proof}

英文文献把理想中每个元素各自幂零的条件称为 nil ideal，把存在共同指数使 \(I^n=0\) 的条件称为 nilpotent ideal。本书将“幂零理想”用于后者；前者不要求各元素具有同一个幂零指数。

\ProseParagraphBreak
对域 \(k\) 和 \(n\ge1\)，在 \(k[t]/(t^n)\) 中，理想 \((\bar t)\) 幂零，故包含于根；商为域，\((\bar t)\) 又是极大理想，所以根恰为它。一般环的根未必幂零，例如 \(k[\![t]\!]\) 的单位恰为常数项非零的级数：逐次解系数可构造其逆。因此非单位元组成的理想 \((t)\) 是唯一极大理想，\(J\bigl(k[\![t]\!]\bigr)=(t)\)，但 \(t^n\ne0\) 对每个 \(n\) 成立。

\subsection{商环、矩阵环与有限长度判据}
\label{b2:subsec:1401:03}

\begin{proposition}\label{b2:prop:radical-quotient}
设 \(R\) 为含幺环、\(I\) 为其双边理想，\(\pi:R\to R/I\) 为商映射，则 \(\pi\bigl(J(R)\bigr)\subseteq J(R/I)\)。若 \(I\subseteq J(R)\)，则
\[
J(R/I)=J(R)/I.
\]
特别地，\(J\bigl(R/J(R)\bigr)=0\)。
\end{proposition}
\begin{proof}
\(R/I\) 的幺性简单左模通过 \(\pi\) 成为简单 \(R\) 模，并被 \(I\) 零化；反过来，满足 \(IS=0\) 的幺性简单 \(R\) 模 \(S\)，其作用也唯一经过 \(R/I\)。两种作用的子模相同，所以简单性不变。因此取商只保留原来简单模中被 \(I\) 零化的那些测试对象。\(J(R)\) 零化全部简单 \(R\) 模，当然也零化这些对象，由\cref{b2:prop:radical-simple-annihilators}得到包含。若 \(I\subseteq J(R)\)，则每个极大左理想都含 \(I\)；商与原像使它们与 \(R/I\) 的极大左理想一一对应。逐个取交便得等式，再取 \(I=J(R)\) 得最后断言。
\end{proof}

测试简单模的族缩小后，共同零化它们的元素可能增加，所以上述包含可以严格。例如 \(J(\mathbb Z)=0\)：若整数 \(x\ne0\)，则 \(1-3x\) 不是整数单位，单位判据排除 \(x\)。但 \(\mathbb Z/4\mathbb Z\) 的根为非零理想 \((\bar2)\)。

\begin{proposition}\label{b2:prop:matrix-radical}
设 \(R\) 为含幺环、\(n\ge1\) 为整数。每个幺性简单左 \(R\) 模 \(S\) 都给出按矩阵乘列作用的幺性简单左 \(M_n(R)\) 模 \(S^n\)。反过来，对任意幺性简单左 \(M_n(R)\) 模 \(N\)，\(E_{11}N\) 通过角环 \(E_{11}M_n(R)E_{11}\cong R\) 成为幺性简单左 \(R\) 模，且
\[
N\cong(E_{11}N)^n
\]
为左 \(M_n(R)\) 模同构。这两个构造给出两类简单模同构类的一一对应。这里 \(E_{ij}\) 为标准矩阵单位。相应地，有
\[
J\bigl(M_n(R)\bigr)=M_n\bigl(J(R)\bigr).
\]
\end{proposition}
\begin{proof}
先把两种环的简单模对应起来。取幺性简单左 \(R\) 模 \(S\)。任意非零列向量有非零坐标 \(s\)，矩阵单位可把它单独移到任意坐标，而 \(Rs=S\)，所以该向量生成全部 \(S^n\)，证明 \(S^n\) 简单。

反过来，取任意幺性简单左 \(M_n(R)\) 模 \(N\)，令 \(S=E_{11}N\)。角环 \(E_{11}M_n(R)E_{11}\) 由仅第 \((1,1)\) 个条目可能非零的矩阵组成，\(r\mapsto rE_{11}\) 是从 \(R\) 到该角环的含幺环同构，角环的单位为 \(E_{11}\)。因此按 \(r\cdot s=(rE_{11})s\) 定义的左 \(R\) 作用是幺性的。这里 \(rE_{ij}\) 表示第 \((i,j)\) 个条目为 \(r\)、其余条目为零的矩阵。

因为 \(1=\sum_iE_{ii}\)，非零 \(v\in N\) 总有某个 \(E_{ii}v\ne0\)。此时 \(E_{1i}v\ne0\)，否则再乘 \(E_{i1}\) 会得到 \(E_{ii}v=0\)，所以 \(S\ne0\)。若 \(0\ne U\subseteq S\) 是 \(R\) 子模，对任意矩阵 \(B=(b_{ji})\) 有
\[
B E_{i1}u=\sum_j E_{j1}\bigl((b_{ji}E_{11})u\bigr)
\qquad(u\in U).
\]
右端仍属于 \(\sum_jE_{j1}U\)，故这个非零子模在所有矩阵作用下稳定，由 \(N\) 简单得 \(\sum_iE_{i1}U=N\)。再施加 \(E_{11}\)，便有 \(U=S\)，所以 \(S\) 简单。

矩阵单位还给出坐标同构 \(S^n\to N\)、\((s_i)\mapsto\sum_iE_{i1}s_i\)。它的逆取 \(v\mapsto(E_{1i}v)_i\)：\(E_{1i}\) 读取第 \(i\) 个坐标，\(E_{i1}\) 将其放回原位置，由 \(E_{1j}E_{i1}=\delta_{ji}E_{11}\) 与 \(\sum_iE_{i1}E_{1i}=1\) 得到两边互逆。若 \(v=\sum_jE_{j1}s_j\)，则
\[
E_{1i}(Bv)=\sum_j(b_{ij}E_{11})s_j=\sum_j b_{ij}\cdot s_j,
\]
正是矩阵乘列的第 \(i\) 个坐标，所以这是左 \(M_n(R)\) 模同构。对原先由 \(S\) 构造的 \(S^n\)，\(E_{11}S^n\) 就是第一坐标上的 \(S\) 副本；模同构也保持这些构造，所以它们在同构类上互逆。

现在，一个矩阵属于 \(J(M_n(R))\)，当且仅当它零化全部简单模 \(S^n\)。把任意 \(s\in S\) 单独放在第 \(j\) 个坐标，就知这等价于每个条目 \(b_{ij}\) 零化全部简单 \(R\) 模。由\cref{b2:prop:radical-simple-annihilators}，又等价于所有 \(b_{ij}\in J(R)\)，得到根公式。零环的两边都为零，直接成立。
\end{proof}

环的根由极大左理想的交给出，模内全部极大子模的交也记录了所有简单商共同检测不到的部分。当这个交为零时，有限长度使这些简单商足以把整个模嵌入有限个简单模的直和，从而给出半单性；这一步中的有限性条件不能删除。

\begin{proposition}\label{b2:prop:finite-length-radical-zero}
设 \(R\) 为含幺环，\(M\) 为有限长度幺性左 \(R\) 模。若 \(M\) 的全部极大子模之交为零，则 \(M\) 半单。因此若正则左模 \({}_RR\) 有限长度，\(R/J(R)\) 为半单环。
\end{proposition}
\begin{proof}
零模直接成立。非零有限长度模有极大子模，由\cref{b2:thm:finite-length-chain-conditions,b2:prop:chain-maximal-minimal}的升链条件及极大元判据可得。为了把全部极大子模的交压缩为有限交，考察它们的所有非空有限交。降链条件使这一族中有一个极小者 \(K=\bigcap_{i=1}^r L_i\)。若 \(K\ne0\)，取 \(0\ne x\in K\)。全部极大子模交为零，故某个极大子模 \(L\) 不含 \(x\)，于是 \(K\cap L\subsetneq K\)，仍是一个有限交，矛盾。因此已有有限个极大子模满足 \(\bigcap_{i=1}^rL_i=0\)。

映射 \(M\to\bigoplus_{i=1}^rM/L_i\)、\(m\mapsto(m+L_i)_i\) 的核为 \(K=0\)，所以它单射。每个商模简单，目标半单，\cref{b2:prop:semisimple-subquotients}保证其子模也半单，故 \(M\) 半单。

令 \(B=R/J(R)\)。它作为左 \(R\) 模具有有限长度，因为有限长度对取商保持。其 \(R\) 作用已经经过 \(B\)，所以 \(R\) 子模与左 \(B\) 子模相同；这些子模在正则 \(B\) 模中又正是左理想。\cref{b2:prop:radical-quotient}给出 \(J(B)=0\)，故上述判据使它作为 \(R\) 模半单。同样因为两种标量作用的子模相同，简单性及直和分解也相同，\({}_BB\) 半单，亦即 \(B\) 是半单环。
\end{proof}

\begin{corollary}[单位的提升]\label{b2:cor:units-lift-radical}
设 \(R\) 为含幺环，\(I\) 为满足 \(I\subseteq J(R)\) 的双边理想。对任意 \(a\in R\)，\(a\) 是 \(R\) 的单位，当且仅当 \(a+I\) 是 \(R/I\) 的单位。更一般地，对任意整数 \(n\ge1\) 和矩阵 \(A\in M_n(R)\)，\(A\) 可逆，当且仅当逐条目取商所得矩阵 \(\overline A\in M_n(R/I)\) 可逆。
\end{corollary}
\begin{proof}
保持单位元的商映射将单位送到单位，所以只需证明反向。若 \(a+I\) 可逆，取其逆的一个提升 \(b\in R\)，有 \(ab=1+i\)、\(ba=1+j\)，其中 \(i,j\in I\)。由\cref{b2:thm:jacobson-unit-criterion}，\(1+i\) 与 \(1+j\) 都可逆。于是
\[
a\bigl(b(1+i)^{-1}\bigr)=1,
\qquad
\bigl((1+j)^{-1}b\bigr)a=1.
\]
一个左逆与一个右逆必相等：若 \(da=1\)、\(ac=1\)，则 \(d=d(ac)=(da)c=c\)。因此 \(a\) 有双侧逆。

逐条目取商的环同态 \(M_n(R)\to M_n(R/I)\) 满射，核为 \(M_n(I)\)，所以目标同构于 \(M_n(R)/M_n(I)\)。又由\cref{b2:prop:matrix-radical}，有 \(M_n(I)\subseteq M_n(J(R))=J(M_n(R))\)。将刚证的单位提升结论用于这个矩阵环及其理想，便得到矩阵版本。
\end{proof}

\begin{proposition}\label{b2:prop:triangular-radical}
设 \(k\) 为域、\(n\ge1\) 为整数，\(T_n(k)\) 为 \(n\times n\) 上三角矩阵组成的含幺 \(k\) 代数，\(N\) 为其中严格上三角矩阵组成的双边理想。则
\[
J(T_n(k))=N,\qquad N^n=0,\qquad T_n(k)/N\cong k^n.
\]
对 \(n\ge2\)，有 \(N^{n-1}\ne0\)，故 \(N\) 的最小幂零指数恰为 \(n\)；\(n=1\) 时 \(N=0\)。
\end{proposition}
\begin{proof}
取对角条目的映射 \(\delta:T_n(k)\to k^n\) 保持加法、乘法和单位元：上三角矩阵乘积的第 \((i,i)\) 个条目恰为两个第 \((i,i)\) 个条目的乘积。它满射，核正是 \(N\)，因此 \(N\) 是双边理想，且 \(T_n(k)/N\cong k^n\)。

严格上三角矩阵中非零条目只能沿严格递增的行列指标出现。\(n\) 个这样的矩阵相乘时，任一可能非零项都需要 \(n+1\) 个严格递增的指标，而指标总共只有 \(n\) 个，所以乘积为零，\(N^n=0\)。由\cref{b2:cor:nil-ideal-radical}，\(N\subseteq J(T_n(k))\)。\(k^n\) 的各坐标模 \(k\) 都简单，而且共同忠实，所以由\cref{b2:prop:radical-simple-annihilators}，\(J(k^n)=0\)。再用\cref{b2:prop:radical-quotient}，根在对角商中的像为零，故 \(J(T_n(k))\subseteq N\)，得到等式。

若 \(n\ge2\)，标准矩阵单位的乘积 \(E_{12}E_{23}\cdots E_{n-1,n}=E_{1n}\ne0\)，属于 \(N^{n-1}\)，所以没有小于 \(n\) 的幂零指数。\(n=1\) 时严格上三角部分为空，\(N=0\)，所述结论直接成立。
\end{proof}
